\documentclass[12pt, a4paper]{article}

% --- Packages ---
\usepackage[utf8]{inputenc}
\usepackage[T1]{fontenc}
\usepackage{geometry}
\geometry{left=2.5cm, right=2.5cm, top=2.5cm, bottom=2.5cm}
\usepackage{times}
\usepackage{natbib}
\setcitestyle{authoryear,open={(},close={)}}
\usepackage{hyperref}
\usepackage{graphicx}
\usepackage{titlesec}
\usepackage{enumitem}
\usepackage{setspace}
\usepackage{amsmath}
\usepackage{amssymb}
\usepackage{algorithm}
\usepackage{algorithmic}
\usepackage{booktabs}
\usepackage{array}
\usepackage{multirow}
\usepackage{xcolor}

\hypersetup{
    colorlinks=true,
    linkcolor=blue,
    filecolor=magenta,
    urlcolor=cyan,
    citecolor=blue,
    pdfencoding=auto
}

\onehalfspacing
\setlength{\parskip}{0.5em}
\titlespacing{\section}{0pt}{14pt}{6pt}
\titlespacing{\subsection}{0pt}{10pt}{4pt}
\titlespacing{\subsubsection}{0pt}{8pt}{3pt}

% =========================================================
\title{\textbf{Multi-Agent Autonomous Driving: A Structured Survey of\\
Learning, Interaction, and Cognitive Paradigms}}
\author{\large Literature Review}
\date{\today}

\begin{document}

\maketitle

% =========================================================
\begin{abstract}
Multi-agent autonomous driving (MAAD) presents one of the most demanding
challenges in modern artificial intelligence: enabling vehicles to navigate complex,
dynamic traffic environments through coordinated perception, prediction, and
decision-making. This survey provides a systematic review of the field, organized
around a clear three-layer taxonomy. We first establish the formal foundations,
framing single-agent driving as a Partially Observable Markov Decision Process
(POMDP) and extending it to the Partially Observable Stochastic Game (POSG)
model required for principled multi-agent reasoning. We then survey
\textbf{(i) single-agent learning foundations}, covering behavior cloning from raw
pixels \citep{bojarski2016end}, perturbation-based data synthesis to overcome
covariate shift \citep{bansal2018chauffeur}, and policy gradient optimization
\citep{schulman2017proximal}; \textbf{(ii) multi-agent interaction and coordination},
examining the evolution from social-force grid pooling \citep{alahi2016social} and
vectorized graph neural networks \citep{gao2020vectornet} to bi-directional
prediction--planning coupling \citep{song2020pip}, Centralized Training
Decentralized Execution (CTDE) reinforcement learning \citep{lowe2017multi},
safety-constrained hierarchical RL \citep{shalev2016safe}, cooperative
perception via V2X feature sharing \citep{wang2020v2vnet}, and scalable
multi-agent training platforms \citep{palanisamy2020multi, zhou2020smarts}; and
\textbf{(iii) LLM-based cognitive agents}, analyzing how large language models
inject common-sense reasoning, memory-augmented adaptation, and interpretable
decision chains into the driving stack
\citep{fu2023drive, wen2023dilu, mao2023gpt, sha2023languagempc}.
Across all paradigms, we synthesize the key technical contributions, identify the
precise limitations each work addresses, and draw cross-cutting connections.
We conclude with a structured analysis of open challenges: formal safety
verification under distributional shift, bandwidth-constrained cooperative
perception, real-time LLM inference, and the need for generative world models
to cover long-tail scenarios.
\end{abstract}

\noindent\textbf{Keywords:} Multi-Agent Autonomous Driving, Imitation Learning,
Deep Reinforcement Learning, Multi-Agent Reinforcement Learning, Graph Neural
Networks, Trajectory Prediction, Vehicle-to-Vehicle Communication,
Large Language Models, Cognitive Agents.

\newpage
\tableofcontents
\newpage

% ---------------------------------------------------------
\section{Introduction}
\label{sec:intro}

Autonomous driving (AD) is a canonical testbed for intelligent systems research: it
demands real-time fusion of noisy, high-dimensional sensor inputs; prediction of the
intentions of heterogeneous road users; and execution of safe, efficient maneuvers
in a continuously evolving, \emph{multi-agent} environment. Despite decades of
engineering effort, achieving the robustness and generalization required for
unrestricted public deployment remains an open problem.

The historical arc of the field can be partitioned into three successive paradigm
shifts, each motivated by well-defined limitations of its predecessor.

\paragraph{Paradigm 1 -- Modular Pipelines and Early Learning.}
Classical autonomous driving relied on a modular perception $\to$ prediction
$\to$ planning $\to$ control stack. The promise of data-driven learning was first
demonstrated at scale by \citet{bojarski2016end}, who showed that a convolutional
network mapping raw front-camera images directly to steering commands could
navigate real roads. This end-to-end approach exposed a fundamental problem:
\emph{covariate shift}, where test-time states diverge from the training distribution
and recovery behaviors were never learned. \citet{bansal2018chauffeur} addressed
this with systematic perturbation and synthesis of adverse scenarios, a technique
that proved even 30 million expert demonstrations were insufficient without such
augmentation.

\paragraph{Paradigm 2 -- Interaction-Aware Multi-Agent Systems.}
Urban driving is inherently multi-agent: the optimal action for the ego-vehicle
depends on the intentions of surrounding agents, which in turn depend on the
ego-vehicle's own plan. Handling this mutual dependence required two independent
advances: (a) \emph{representation} methods that encode the relational structure of
traffic scenes compactly and efficiently, and (b) \emph{policy learning} algorithms
that remain stable when multiple agents adapt simultaneously. On the representation
front, the field evolved from grid-based social pooling \citep{alahi2016social} to
vectorized graph neural networks \citep{gao2020vectornet}, and further to
bi-directional prediction--planning coupling \citep{song2020pip}. On the policy
learning front, the Centralized Training Decentralized Execution (CTDE) paradigm
of \citet{lowe2017multi}, safety-constrained hierarchical RL \citep{shalev2016safe},
connected multi-agent platforms \citep{palanisamy2020multi}, scalable training
environments \citep{zhou2020smarts}, and intermediate-feature V2X sharing
\citep{wang2020v2vnet} collectively defined the state of the art.

\paragraph{Paradigm 3 -- LLM-Based Cognitive Agents.}
Even well-trained multi-agent systems remain fragile in \emph{long-tail} scenarios
(edge cases outside the statistical support of the training distribution). The
emergence of large language models (LLMs) with broad world knowledge, chain-of-thought
reasoning, and few-shot generalization has inspired a new class of \emph{cognitive
driving agents}. \citet{fu2023drive} demonstrated that LLMs can reason through
novel scenes step-by-step; \citet{wen2023dilu} equipped agents with
reflection-and-memory mechanisms for lifelong adaptation; \citet{mao2023gpt}
reformulated motion planning as a language modeling task; and
\citet{sha2023languagempc} introduced a hierarchical neuro-symbolic architecture
pairing LLM strategic reasoning with MPC execution.

\subsection{Scope and Contributions}
This survey covers 18 primary papers spanning 2016--2025. Our contributions over
existing reviews \citep{kiran2021deep} are threefold. First, we adopt an explicit
multi-agent framing throughout, placing all methods within either the POMDP or
POSG formalism. Second, we construct a coherent three-layer taxonomy that
distinguishes foundational single-agent learning from multi-agent interaction and
from cognitive-level reasoning, enabling systematic comparison. Third, we dedicate
a structured section to LLM-based paradigms, which have emerged predominantly
after existing surveys were written. For each method we analyze: (1) the precise
limitation it addresses, (2) the technical mechanism it introduces, and (3) the
residual limitations that motivate subsequent work.

\paragraph{Reference Weighting Strategy.}
To keep the survey concise while preserving depth, we assign \emph{unequal}
coverage to references. Tier-1 landmark works (e.g., DAVE-2, ChauffeurNet,
MADDPG/CTDE, VectorNet, PiP, V2VNet, and representative LLM agents) receive
full technical discussion with equations or architectural analysis. Tier-2
works are used for focused comparison and evidence. Tier-3 works are retained
mainly for context and historical continuity through brief mentions. This
weighted treatment avoids giving all papers equal narrative space and keeps the
overall manuscript within a standard survey length.

\subsection{Paper Organization}
Section~\ref{sec:background} establishes the formal problem framework.
Section~\ref{sec:single} surveys single-agent learning foundations.
Section~\ref{sec:interaction} examines multi-agent interaction and coordination.
Section~\ref{sec:cognitive} reviews LLM-based cognitive agents.
Section~\ref{sec:benchmarks} surveys simulation environments and evaluation
benchmarks.
Section~\ref{sec:challenges} synthesizes open challenges and future directions.
Section~\ref{sec:conclusion} concludes.

% ---------------------------------------------------------
\section{Background and Problem Formulation}
\label{sec:background}

\subsection{Single-Agent Decision-Making: POMDP}
The canonical framework for single-agent autonomous driving is the \emph{Partially
Observable Markov Decision Process} (POMDP), defined by the tuple
$\langle \mathcal{S}, \mathcal{A}, \mathcal{O}, T, R, \Omega, \gamma \rangle$:

\begin{itemize}[noitemsep]
  \item $\mathcal{S}$: world state (ego kinematics, road topology, hidden intents
        of other agents).
  \item $\mathcal{A}$: action space, either discrete (lane-change decisions) or
        continuous ($\delta \in$ steering, $a \in$ throttle/brake).
  \item $\mathcal{O}$: observation space from onboard sensors (camera, LiDAR,
        radar, GPS).
  \item $T(s_{t+1} | s_t, a_t)$: stochastic transition dynamics encoding vehicle
        kinematics and environment evolution.
  \item $R(s_t, a_t)$: reward encoding safety (collision avoidance), efficiency
        (progress toward goal), and comfort (minimizing jerk).
  \item $\Omega(o_t | s_t)$: observation emission model.
  \item $\gamma \in [0, 1)$: discount factor.
\end{itemize}

The goal is to find a policy $\pi(a_t | o_{1:t})$ that maximizes the expected
discounted return:
\begin{equation}
  \pi^* = \arg\max_\pi \; \mathbb{E}_{\tau \sim \pi}
          \!\left[ \sum_{t=0}^{T} \gamma^t R(s_t, a_t) \right].
  \label{eq:pomdp}
\end{equation}

\subsection{Multi-Agent Extension: POSG}
Standard POMDP formulations become inadequate in traffic because the transition
function $T$ depends not only on the ego-vehicle's action $a^0_t$ but also on
the actions of $N-1$ other agents $a^1_t, \dots, a^{N-1}_t$, whose policies are
simultaneously evolving during training. This \emph{non-stationarity} violates the
Markov property from the perspective of any single learner.

\citet{palanisamy2020multi} formally adopts the \emph{Partially Observable
Stochastic Game} (POSG) $\langle \mathcal{N}, \mathcal{S}, \{\mathcal{A}^i\},
\{\mathcal{O}^i\}, T, \{R^i\}, \gamma \rangle$, where $\mathcal{N} =
\{1,\dots,N\}$ is the set of agents and each agent $i$ has its own action space
$\mathcal{A}^i$, observation space $\mathcal{O}^i$, and reward function $R^i$.
The joint transition is:
\begin{equation}
  T(s_{t+1} | s_t, a^1_t, \dots, a^N_t).
\end{equation}
When $R^i = R^j$ for all $i, j$, the game is \emph{fully cooperative}; when
rewards are adversarial, it is \emph{zero-sum}; traffic typically falls into
\emph{mixed} cooperative-competitive settings.

The POSG framing clarifies the two principal challenges of MAAD:
\begin{enumerate}[noitemsep]
  \item \textbf{Non-stationarity}: Other agents adapt, so the environment is
        non-Markovian from any single agent's perspective.
  \item \textbf{Exponential joint space}: The joint action space
        $|\mathcal{A}|^N$ scales exponentially with the number of agents.
\end{enumerate}

\citet{shalev2016safe} further observe that, for policy gradient methods,
the gradient estimator variance grows proportionally to $1/p^2$ for an accident
probability $p \ll 1$ (derived analytically from the REINFORCE variance bounds
in their Lemma~3), making unconstrained RL impractical for safety-critical
applications without architectural safeguards.

\subsection{Architectural Dichotomy: Modular vs.\ End-to-End}
The POMDP/POSG formalism can be instantiated through two broad architectural
choices:

\paragraph{Modular Pipelines.}
Perception $\to$ Prediction $\to$ Planning $\to$ Control subsystems are designed
and validated independently. The advantages are interpretability and modular
testing \citep{caesar2021nuplan}. The critical disadvantage is \emph{cascading
error}: bounding-box representations lose contextual information (e.g., pedestrian
gaze), and the prediction module may produce worst-case forecasts that paralyze
the planner in a well-known ``frozen robot'' failure mode \citep{song2020pip}.

\paragraph{End-to-End Learning.}
A differentiable function $f_\theta: \mathcal{O} \to \mathcal{A}$ is learned
jointly, optimizing the terminal driving objective without hand-crafted intermediate
representations \citep{bojarski2016end}. The entire feature hierarchy is learned
for the task, but the resulting black-box model is hard to debug, certify, or
bound from safety constraints.

The papers reviewed in this survey represent the full spectrum of these
architectural choices, increasingly augmented by LLM-based reasoning components.

% ---------------------------------------------------------
\section{Single-Agent Learning Foundations}
\label{sec:single}

Before multi-agent interaction can be addressed, robust single-agent policies must
first be established. This section examines the foundational works on behavior
cloning and policy optimization that underpin all subsequent MAAD methods.

\subsection{End-to-End Behavior Cloning: DAVE-2}
\label{sec:dave}

\citet{bojarski2016end} from NVIDIA introduced \emph{DAVE-2}, the first large-scale
demonstration of end-to-end convolutional neural network (CNN) driving. A
nine-layer CNN with approximately 27 million connections maps three forward-facing
camera images (with synthetic viewpoint perturbations) directly to steering angle.
The network was trained on 72 hours of human driving using the mean-squared error
behavioral cloning (BC) loss:
\begin{equation}
  \mathcal{L}_{\text{BC}}(\theta) =
    \mathbb{E}_{(o, a) \sim \mathcal{D}}\!\left[ \|\pi_\theta(o) - a \|^2 \right].
  \label{eq:bc}
\end{equation}
A key insight was the use of \emph{synthetic viewpoint perturbation}: by
offsetting the camera laterally and recording the corresponding corrective steering
command, the authors generated recovery behaviors in a computationally cheap way.

\paragraph{Technical Significance.}
DAVE-2 established the feasibility of implicit feature learning for driving (lane
markings, road boundaries, and even distant curve curvature emerged without
explicit supervision labels). However, its failure mode is well-defined:
Equation~\eqref{eq:bc} minimizes one-step prediction error under the training
distribution $\mathcal{D}$, but by a standard BC analysis, this agent
accrues \emph{quadratic} compounding error with trajectory length because the
agent, once displaced, has no recovery mechanism. This covariate shift problem
directly motivated the work of \citet{bansal2018chauffeur}.

\subsection{Overcoming Covariate Shift: ChauffeurNet}
\label{sec:chauffeur}

\citet{bansal2018chauffeur} demonstrated that even a dataset of 30 million
human-driving clips fails to prevent covariate shift under standard BC. Their key
insight is that an agent trained only on \emph{optimal} demonstrations never
observes perturbed or adversarial states, so its policy degrades catastrophically
once it drifts from nominal conditions.

ChauffeurNet addresses this through \emph{imitation with synthesized perturbation}.
The system processes rendered mid-level scene representations (agent bounding
boxes, road map, speed, traffic lights) rather than raw pixels, enabling controlled
numerical augmentation. Training adds:
\begin{enumerate}[noitemsep]
  \item \textbf{Geometric perturbation}: Expert trajectories are warped
        randomly, injecting off-route and near-collision states.
  \item \textbf{Environment dropout}: Randomly masking objects in the input forces
        the agent to learn robust representations.
  \item \textbf{Loss on collision and off-road events}: Auxiliary losses penalize
        synthesized near-accident trajectories and encourage active recovery.
\end{enumerate}
The aggregate loss is
\begin{equation}
  \mathcal{L} = \mathcal{L}_{\text{imitation}}
              + \lambda_c\,\mathcal{L}_{\text{collision}}
              + \lambda_r\,\mathcal{L}_{\text{road}}
              + \lambda_g\,\mathcal{L}_{\text{geometry}},
\end{equation}
where $\lambda_c, \lambda_r, \lambda_g$ weight the collision, road-geometry, and
trajectory-geometry penalties respectively. This \emph{synthesize the worst}
strategy expands the training distribution to cover recovery states without
collecting real data for every possible perturbation.

\paragraph{Residual Limitation.}
ChauffeurNet still treats other agents as passive objects in the scene
representation; it does not explicitly model their future intentions or strategic
interactions. This limitation is addressed by the multi-agent interaction methods
of Section~\ref{sec:interaction}.

\subsection{Policy Optimization: PPO}
\label{sec:ppo}

When closed-loop interaction is required, RL methods can directly optimize
Equation~\eqref{eq:pomdp}. The most widely deployed policy gradient algorithm in
the AD community is \emph{Proximal Policy Optimization} (PPO)
\citep{schulman2017proximal}. PPO replaces the second-order trust-region
constraint of TRPO with a first-order \emph{clipped surrogate objective}:
\begin{equation}
  L^{\text{CLIP}}(\theta)
  = \mathbb{E}_t \!\left[
      \min\!\left( r_t(\theta)\hat{A}_t,\;
            \text{clip}(r_t(\theta), 1-\epsilon, 1+\epsilon)\hat{A}_t
      \right)
    \right],
  \label{eq:ppo}
\end{equation}
where $r_t(\theta) = \pi_\theta(a_t|s_t)/\pi_{\theta_{\text{old}}}(a_t|s_t)$ is
the probability ratio, $\hat{A}_t$ is the generalized advantage estimate, and
$\epsilon$ (typically 0.1--0.2) clips the ratio to prevent excessively large
policy updates. The clipping removes the incentive for the optimizer to move the
policy far from $\pi_{\theta_{\text{old}}}$ in a single step, providing
monotonic improvement guarantees without solving a quadratic programming subproblem.

\paragraph{Relevance to MAAD.}
PPO has been adopted as the default optimizer in multi-agent training loops
(MACAD-Gym \citep{palanisamy2020multi}, SMARTS \citep{zhou2020smarts}) because
its conservative update rule reduces the risk of policy collapse when the apparent
reward landscape shifts as co-agent policies evolve. \citet{kiran2021deep} provide
a comprehensive survey situating PPO and related DRL methods (DQN, DDPG, SAC)
within the full spectrum of autonomous driving tasks.

% ---------------------------------------------------------
\section{Multi-Agent Interaction and Coordination}
\label{sec:interaction}

Having established single-agent foundations, this section examines how the field
has addressed the two POSG challenges identified in Section~\ref{sec:background}:
\emph{representing} coupled agent behavior and \emph{training} stable policies in
a non-stationary multi-agent environment.

\subsection{Representing Agent Interactions: From Grid Pooling to Graph Networks}

\subsubsection{Social LSTM: Spatiotemporal Social Pooling}
\label{sec:social}

The earliest learning-based approach to agent interaction was introduced by
\citet{alahi2016social} for pedestrian trajectory prediction. Their insight is that
pedestrians follow implicit social norms (collision avoidance, group cohesion,
lane-keeping on pavements) that cannot be captured by independent Kalman filters.

\emph{Social LSTM} uses a separate LSTM for each agent, sharing hidden states
through a novel \emph{Social Pooling} layer. For agent $i$ at time $t$, the pooled
state is
\begin{equation}
  H_i^{\text{pool}} =
    \sum_{j \in \mathcal{N}_i}
    \mathcal{E}\!\left( \mathbb{1}_{[x_j, y_j \in \text{cell}(m,n)]} \cdot
                  W_e\, h_{j,t-1} \right),
\end{equation}
where $\mathcal{N}_i$ are neighbors within a spatial threshold, the Boolean
indicator selects neighbors in occupancy-grid cell $(m,n)$, $W_e$ is a learned
embedding, and $\mathcal{E}$ is an element-wise max operation. The pooled tensor
is concatenated with the agent's own embedding and fed as input to its LSTM.

\paragraph{Significance and Limitations.}
Social LSTM demonstrated that trajectory prediction is fundamentally a
\emph{relational} problem; ignoring neighbor states leads to physically
implausible predictions (e.g., interpenetrating bounding boxes). However, the
fixed-size occupancy grid introduces two structural limitations: (1) spatial
resolution is a hyperparameter that cannot adapt to varying agent densities, and
(2) representing a large intersection or motorway segment requires a very large
grid, making computation prohibitive. These limitations motivated the vectorized
graph approach of VectorNet.

\subsubsection{VectorNet: Vectorized Scene Representation with GNNs}
\label{sec:vectornet}

\citet{gao2020vectornet} introduced a paradigm shift from rasterization to
\emph{vectorization}. Instead of rendering the scene into a pixel grid, VectorNet
represents every map element (lane centerline, crosswalk boundary, stop line) and
every agent's historical trajectory as an ordered sequence of vectors (polylines).
Each vector $v_i$ encodes start point, end point, and semantic attributes.

VectorNet employs a two-level hierarchical GNN:
\begin{enumerate}[noitemsep]
  \item \textbf{Subgraph network}: Within each polyline (e.g., one lane), a local
        GNN encodes pairwise vector relations into a polyline-level node feature
        $p_i = \text{GNN}_{\text{local}}(\{v_j : v_j \in \text{polyline}_i\})$.
  \item \textbf{Global interaction graph}: All polyline nodes $\{p_i\}$ form a
        fully connected graph. Multi-head self-attention models global
        interdependencies between lane topology and dynamic agents:
        \begin{equation}
          p_i' = \text{Attention}\!\left( p_i W_Q,\; \{p_j\} W_K,\; \{p_j\} W_V \right).
          \label{eq:vn_attn}
        \end{equation}
\end{enumerate}
The vectorized representation reduces parameter count by over 70\% compared to
CNN-based renderers and achieves state-of-the-art performance on Argoverse motion
forecasting, because the GNN can directly exploit structural adjacency (e.g., a
lane successor relation) that would be invisible to a convolutional filter.

\paragraph{Relevance to Multi-Agent Settings.}
VectorNet's global attention (Equation~\ref{eq:vn_attn}) scales to arbitrarily
many agents and map elements without grid-size constraints, making it the
predominant scene encoder in modern MAAD stacks. The learned attention weights
also provide interpretable evidence of which agents and lanes most influenced
each prediction.

\subsubsection{PiP: Closing the Prediction--Planning Loop}
\label{sec:pip}

A structural limitation shared by all methods described so far is that prediction
and planning are \emph{decoupled}: the predictor forecasts other agents'
trajectories independently of the ego-vehicle's future actions. This creates the
``frozen robot problem'' identified by \citet{song2020pip}: if the ego-vehicle's
intent is unknown, a conservative predictor assigns high probability to trajectories
that cut off the ego-vehicle, causing the planner to halt indefinitely.

\citet{song2020pip} addressed this with \emph{Planning-informed Prediction} (PiP).
The key innovation is to \emph{condition} the prediction module on a set of
candidate ego-vehicle trajectories $\{\tau^k_{\text{ego}}\}_{k=1}^K$:
\begin{equation}
  \hat{Y}_i = f_{\text{pred}}\!\left( X_i,\; X_{-i},\; \tau^k_{\text{ego}} \right),
  \quad k \in \{1, \dots, K\},
\end{equation}
where $X_i$ is the past trajectory of target agent $i$, $X_{-i}$ are other
contextual agents, and $\tau^k_{\text{ego}}$ informs the predictor that the
ego-vehicle \emph{intends} a specific maneuver (e.g., lane merge). The planning
module then selects the ego trajectory whose induced prediction is globally safest
and most efficient.

PiP demonstrated on nuScenes that this feedback loop reduces collision rate and
improves trajectory realism, confirming the game-theoretic view that driving is a
\emph{bi-directional} interaction: the ego-vehicle's plan shapes others' reactions,
which in turn constrains the optimal plan.

\subsection{Multi-Agent Reinforcement Learning}

\subsubsection{Safety-Constrained Hierarchical RL}
\label{sec:rss}

\citet{shalev2016safe} identify two fundamental obstacles to applying RL directly
to autonomous driving: (1) the non-stationarity of the multi-agent environment
invalidates the Markov assumption, and (2) the gradient variance of policy gradient
methods scales as $O(1/p^2)$ for accident probability $p$, making learning
impractically slow when safety is enforced only through a large negative reward.

Their solution decomposes the policy into two components:
\begin{enumerate}[noitemsep]
  \item \textbf{Policy for Desires ($\pi_D$)}: A deep neural network, trained by
        policy gradient, encodes preferences for comfort and efficiency (smooth
        speed, comfortable turns). This component is \emph{learned}.
  \item \textbf{Hard Safety Constraints}: A trajectory planner enforces inviolable
        physical and legal constraints (safe following distance, right-of-way) as
        hard bounds on the actions proposed by $\pi_D$. This component is
        \emph{not learned}.
\end{enumerate}
The resulting \emph{Responsibility-Sensitive Safety} (RSS) architecture acts as a
hierarchical filter: RL explores freely within a provably safe action subset,
circumventing the variance problem because catastrophic-reward trajectories are
never executed.

Crucially, \citet{shalev2016safe} also show that \emph{policy gradient does not
require the Markov assumption}: the REINFORCE likelihood-ratio gradient estimator
remains unbiased when the environment transition depends on other agents, as long
as their policies are fixed (or slowly varying) during a rollout. This observation
provides theoretical justification for applying policy gradient in MAAD settings.

\subsubsection{MADDPG: Centralized Training, Decentralized Execution}
\label{sec:maddpg}

\citet{lowe2017multi} address the non-stationarity problem at the algorithmic level.
Naive Q-learning fails in multi-agent settings because the non-stationary returns
invalidate the Bellman optimality condition. Policy gradient methods accumulate
variance proportional to the number of agents, since each agent's gradient is
computed under uncertainty about co-agents' actions.

The \emph{Multi-Agent Deep Deterministic Policy Gradient} (MADDPG) algorithm
resolves this by adopting a \emph{centralized critic, decentralized actor} design:
\begin{itemize}[noitemsep]
  \item \textbf{Centralized critic}: $Q_i(\mathbf{o}, a^1, \dots, a^N)$ takes the
        joint observation and action tuple as input during training, producing a
        stable value function because the joint environment \emph{is} stationary
        conditioned on all agents' actions.
  \item \textbf{Decentralized actor}: $\pi_i(a^i | o^i)$ conditions only on local
        observation $o^i$ at execution time, enabling deployment without
        infrastructure-level coordination.
\end{itemize}
To handle cooperative-competitive mixed settings, the actor is updated via the
centralized critic's gradient:
\begin{equation}
  \nabla_{\theta_i} J(\theta_i)
  = \mathbb{E}\!\left[
      \nabla_{\theta_i} \log \pi_i(a^i|o^i)\;
      Q_i^\pi(\mathbf{o}, a^1, \dots, a^N)
    \right].
  \label{eq:maddpg}
\end{equation}
An ensemble of policies is further maintained per agent to encourage robustness to
diverse partner strategies.

\paragraph{Impact on MAAD.}
The CTDE paradigm has become the dominant paradigm for MAAD: SMARTS
\citep{zhou2020smarts}, MACAD-Gym \citep{palanisamy2020multi}, and most subsequent
works adopt centralized critics with decentralized execution. The framework cleanly
separates the concerns of training infrastructure (which may have access to global
state via V2X or roadside units) from real-time deployment (which must act on local
observations alone).

\subsubsection{MACAD-Gym: Connected Autonomous Driving Platform}
\label{sec:macad}

\citet{palanisamy2020multi} operationalize the POSG formulation into an open
research platform, \emph{MACAD-Gym}, built on top of the CARLA simulator
\citep{dosovitskiy2017carlaopenurbandriving}. The platform provides:
\begin{itemize}[noitemsep]
  \item A set of multi-agent connected AD scenarios (urban intersections, ramps,
        merges) with heterogeneous agent types.
  \item A taxonomy of multi-agent environments along three axes: (a) nature of
        tasks (cooperative/competitive/mixed), (b) nature of agents
        (homogeneous/heterogeneous), and (c) nature of the environment
        (stationary/non-stationary).
  \item An extensible gym-compatible API enabling plug-in of any RL algorithm.
\end{itemize}
Baseline MACAD-Agents trained with PPO on raw camera observations demonstrate
successful policy learning in a partially observable stop-sign-controlled
three-way intersection. The platform's value lies in making the abstract POSG
formalism concrete and reproducible.

\subsubsection{SMARTS: Scalable Multi-Agent RL Training School}
\label{sec:smarts}

\citet{zhou2020smarts} identify a practical gap: existing simulators either lack
realistic multi-agent traffic behavior models or lack the programmatic API needed
for MARL research. \emph{SMARTS} addresses this with three design principles:
\begin{enumerate}[noitemsep]
  \item \textbf{Diverse social vehicle behavior}: Realistic background vehicles are
        modeled with behavior models including IDM (car-following), MOBIL
        (lane-change), and custom aggressive/distracted archetypes, avoiding the
        problem of agents that learn to exploit unrealistic passive opponents.
  \item \textbf{Ego-agent heterogeneity}: Multiple ego-agent types (sedan,
        motorcycle, bus) with distinct dynamics can be trained simultaneously.
  \item \textbf{Reproducible scenario definitions}: Scenarios are declaratively
        specified in a human-readable format, ensuring reproducibility across labs.
\end{enumerate}
SMARTS introduced a large-scale empirical analysis showing that social vehicle
diversity substantially affects the generalization of trained policies, highlighting
a benchmark-quality issue in prior work where opponents were too predictable.

\subsection{Cooperative Perception via V2X: V2VNet}
\label{sec:v2vnet}

All methods discussed so far rely on \emph{onboard} sensing only. \citet{wang2020v2vnet}
argue that this imposes a fundamental ceiling on safety: a single viewpoint cannot
see around occlusions (a child chasing a ball behind a parked van, for instance).
Vehicle-to-Vehicle (V2V) communication can, in principle, share observations from
multiple viewpoints; the practical challenge is choosing \emph{what} to share.

\citet{wang2020v2vnet} analyze three candidates:
\begin{enumerate}[noitemsep]
  \item \textbf{Raw sensor data} (LiDAR point clouds): High bandwidth requirement;
        infeasible within the V2X bandwidth constraints analyzed by \citet{wang2020v2vnet}.
  \item \textbf{High-level tracks} (3D bounding boxes): Low bandwidth but loses
        context about object appearance and uncertainty.
  \item \textbf{Intermediate feature maps}: Compressed activations from an
        intermediate backbone layer strike the optimal bandwidth--accuracy tradeoff.
\end{enumerate}
V2VNet implements option 3 by having each SDV broadcast its compressed feature map
$\mathbf{F}_j^{\text{ego}}$. The receiving vehicle aggregates features from all
$j \in \mathcal{N}(i)$ (within 70\,m) using a spatially-aware GNN:
\begin{equation}
  \mathbf{F}_i^{\text{fused}}
  = \text{GNN}_{\text{agg}}\!\left(
      \mathbf{F}_i^{\text{ego}},\;
      \{T_{j \to i}(\mathbf{F}_j^{\text{ego}})\}_{j \in \mathcal{N}(i)}
    \right),
\end{equation}
where $T_{j \to i}$ applies a learned spatial transformation to align vehicle
$j$'s feature map into vehicle $i$'s coordinate frame. The fused feature map
feeds a joint perception and prediction (P\&P) head.

Experiments on the authors' V2VSim dataset demonstrate that V2VNet substantially
improves detection AP and motion forecasting ADE/FDE in scenarios with occlusions,
while meeting bandwidth constraints. This work opens the direction of
\emph{cooperative MAAD}: agents that not only plan cooperatively but also
\emph{perceive} cooperatively.

% ---------------------------------------------------------
\section{LLM-Based Cognitive Agents}
\label{sec:cognitive}

The methods surveyed in Sections~\ref{sec:single}--\ref{sec:interaction} are all
trained by optimization: they minimize empirical risk or maximize expected reward
over a fixed training distribution. Their shared Achilles' heel is the
\emph{long-tail problem}: rare, semantically complex scenarios (a funeral
procession blocking a lane; a police officer overriding a green light) are
statistically underrepresented in any finite training corpus, yet carry
disproportionate safety significance.

Large language models (LLMs) offer a qualitatively different resource: broad world
knowledge acquired from internet-scale text, the ability to perform multi-step
chain-of-thought reasoning, and few-shot adaptation from task descriptions rather
than gradient updates. This section surveys the four most influential attempts to
harness these capabilities for autonomous driving.

\subsection{Drive Like a Human: Chain-of-Thought Scene Reasoning}
\label{sec:dlh}

\citet{fu2023drive} argue that the defining difference between a novice and an
expert human driver is not reflexes but \emph{contextual reasoning}: the ability
to interpret an ambiguous situation (a cyclist drifting toward the lane center)
as having a probable cause (avoiding a pothole) and responding appropriately.
They frame this as a question: can an LLM replicate this reasoning from natural
language scene descriptions?

Their system employs \emph{Chain-of-Thought} (CoT) prompting, structuring the LLM
response as a three-stage inference trace:
\begin{enumerate}[noitemsep]
  \item \textbf{Observation}: Translate the sensor-derived scene state into a
        natural language description (``A pedestrian is standing at the crosswalk
        edge, looking left'').
  \item \textbf{Reasoning}: The LLM infers probable intent
        (``They are checking for oncoming traffic; they will likely cross shortly'').
  \item \textbf{Decision}: The LLM generates an actionable response
        (``Reduce speed to 15\,km/h and prepare to yield'').
\end{enumerate}
Empirical evaluation on a curated set of long-tail scenarios demonstrates
substantially higher success rates than purely end-to-end models, particularly for
scenarios involving implicit social context and non-standard road configurations.

\paragraph{Limitations.}
The system outputs textual decisions, which must be translated to low-level
controls by a downstream interpreter. It does not incorporate a memory of past
failures, so the same reasoning error can recur. Both limitations are addressed by
DiLu (Section~\ref{sec:dilu}) and LanguageMPC (Section~\ref{sec:lmpc}).

\subsection{DiLu: Memory-Augmented Knowledge-Driven Driving}
\label{sec:dilu}

\citet{wen2023dilu} propose that an ideal driving agent should behave like an
experienced driver: one who not only applies general knowledge to new situations
but also recalls and applies lessons from previous mistakes. They introduce
\emph{DiLu}, a four-component framework:

\begin{enumerate}
  \item \textbf{Environment Interaction Module}: A vision-language encoder
        translates the visual scene into a structured textual description, including
        agent identities, relative positions, velocities, and traffic signals.

  \item \textbf{Reasoning Module}: An LLM (GPT-4) receives the scene description
        and retrieved memories, then generates a decision with an explicit
        chain-of-thought rationale.

  \item \textbf{Reflection Module}: After execution, if a safety criterion is
        violated or the reward falls below a threshold, the reflection module
        prompts the LLM to perform root-cause analysis (``The rear vehicle braked
        harder than expected when I changed lanes; I underestimated its speed'').
        This analysis is encoded as a \emph{lesson}.

  \item \textbf{Memory Bank}: Lessons are stored as dense vector embeddings using
        a sentence encoder. At inference time, the $k$ nearest lessons (by cosine
        similarity) are retrieved and prepended to the LLM prompt, enabling the
        reasoning module to avoid previously encountered failure modes.
\end{enumerate}

The retrieval-augmented generation loop enables \emph{few-shot continual learning}:
DiLu improves from episode to episode without gradient updates to the LLM backbone,
a form of meta-adaptation critical for deployment in diverse and changing
environments. Experiments in CARLA and a closed-loop highway simulator show that
DiLu's reflection-and-memory mechanism substantially outperforms the same LLM
backbone used without memory.

\paragraph{Distinction from Prior Work.}
Unlike end-to-end models that encode experience implicitly in network weights,
DiLu externalizes experience into an interpretable, queryable memory. Unlike
Drive Like a Human, DiLu closes the loop from decision to outcome, enabling
active error correction. The limitation is that the quality of the reflection is
bounded by the LLM's ability to diagnose its own failures---a property not yet
formally verified.

\subsection{GPT-Driver: Motion Planning as Language Modeling}
\label{sec:gptdriver}

\citet{mao2023gpt} take a different angle: rather than using LLMs for high-level
reasoning only, they reformulate \emph{trajectory prediction itself} as a language
generation task. Continuous waypoint coordinates are discretized and tokenized
into a text sequence, and GPT-3.5 is fine-tuned to complete this sequence
conditioned on a scene description:
\begin{equation}
  P(T | O) = \prod_{i=1}^{L} P(w_i | w_{1:i-1}, O),
  \label{eq:gpt_driver}
\end{equation}
where $T = (w_1, \dots, w_L)$ is the tokenized trajectory and $O$ is the
natural language scene observation. Crucially, the scene description also
includes \emph{textual reasoning}: the model first generates a reasoning trace
(``The lead vehicle is decelerating; I should maintain a larger gap'') and then
generates the numerical trajectory conditioned on that trace.

Evaluated on the nuScenes benchmark, GPT-Driver achieves state-of-the-art
displacement errors with fewer parameters than specialized motion planners.
The results suggest that the LLM's sequence modeling capacity, trained on vast
amounts of structured text, generalizes effectively to the structured token
sequences of driving trajectories. Furthermore, the reasoning trace provides
intrinsic interpretability: an operator can read \emph{why} the vehicle chose a
particular path.

\paragraph{Key Insight.}
The reformulation as language modeling enables \emph{in-context learning}: by
providing a few demonstrations of novel scenarios in the prompt, the model
generalizes without fine-tuning. This is qualitatively different from all
prior methods that require gradient updates to adapt to new conditions.

\subsection{LanguageMPC: Neuro-Symbolic Hierarchical Control}
\label{sec:lmpc}

A practical barrier to LLM-based driving is latency: state-of-the-art LLMs
require hundreds of milliseconds per query, whereas safety-critical control
loops demand decisions at 10--20\,Hz. \citet{sha2023languagempc} resolve this
tension through a hierarchical neuro-symbolic architecture.

The system operates on two time scales:
\begin{enumerate}
  \item \textbf{Strategic level (LLM, $\sim$1--2\,Hz)}: The LLM receives a
        natural language scene description and generates a high-level behavioral
        directive (e.g., ``Overtake the slow vehicle by merging left''), together
        with a set of numerical parameters that govern the lower-level objective:
        target lane $l^*$, target speed $v^*$, and weights $\mathbf{w}$ for the
        MPC cost function.
  \item \textbf{Execution level (MPC, 10--20\,Hz)}: A Model Predictive Control
        loop solves a constrained trajectory optimization:
        \begin{equation}
          \min_{u_{0:H}}
          \sum_{k=0}^{H-1} \mathbf{w}^\top \phi(x_k, u_k, l^*, v^*)
          \quad \text{s.t.} \quad x_{k+1} = f(x_k, u_k),\;
          g(x_k, u_k) \leq 0,
          \label{eq:mpc}
        \end{equation}
        where $\phi$ is the feature vector of the cost function, $f$ is the
        vehicle dynamics model, and $g$ encodes hard safety constraints. The MPC
        runs in real time and handles dynamic obstacle avoidance and constraint
        satisfaction that the LLM need not reason about explicitly.
\end{enumerate}

The decoupling is bidirectional: the MPC handles real-time safety, while the LLM
provides semantic adaptability (e.g., adjusting the weight on $v^*$ when passing
a school zone). This ``thinking fast and slow'' architecture directly addresses
the latency and safety-guarantee limitations of purely LLM-based approaches.

\paragraph{Comparison Across LLM Methods.}
Table~\ref{tab:llm_comparison} summarizes the four LLM-based cognitive agents.

\begin{table}[htbp]
\centering
\caption{Comparison of LLM-based cognitive driving agents.}
\label{tab:llm_comparison}
\resizebox{\linewidth}{!}{%
\begin{tabular}{lllll}
\toprule
\textbf{Method} & \textbf{LLM Role} & \textbf{Memory} &
\textbf{Low-Level Ctrl.} & \textbf{Key Strength} \\
\midrule
Drive Like a Human \citep{fu2023drive} & Scene reasoning & None &
External & CoT interpretability \\
DiLu \citep{wen2023dilu} & Reasoning + reflection & Vector store &
External & Continual adaptation \\
GPT-Driver \citep{mao2023gpt} & Trajectory generation & None &
Self (language) & In-context planning \\
LanguageMPC \citep{sha2023languagempc} & Strategic planning & None &
MPC (hard) & Real-time safety \\
\bottomrule
\end{tabular}%
}
\end{table}

% ---------------------------------------------------------
\section{Simulation Environments and Benchmarks}
\label{sec:benchmarks}

Rigorous empirical evaluation is indispensable for MAAD research. This section
surveys the major simulation platforms and evaluation benchmarks, distinguishing
\emph{open-loop} (replay-based) and \emph{closed-loop} (reactive agent) protocols.

\subsection{CARLA: High-Fidelity Urban Simulation}

\citet{dosovitskiy2017carlaopenurbandriving} introduced \emph{CARLA}, an open-source urban driving simulator, which
has become the de facto standard urban driving simulator in the research community.
CARLA provides:
\begin{itemize}[noitemsep]
  \item Physically realistic rendering with configurable weather, lighting, and
        road surface conditions.
  \item A flexible sensor suite: RGB cameras, depth cameras, semantic segmentation
        cameras, LiDAR, radar, GPS, and IMU.
  \item A Python API for programmatic scenario definition and agent interaction.
  \item OpenDRIVE-based HD map support for accurate road geometry.
\end{itemize}
CARLA's modular architecture supports both \emph{autopilot-controlled} and
\emph{algorithm-controlled} agents, enabling controlled experiments where only
the ego-vehicle policy is under test while traffic is managed by a rule-based
autopilot. DiLu \citep{wen2023dilu}, MACAD-Gym \citep{palanisamy2020multi}, and
GPT-Driver ablations all use CARLA as their primary environment.

\subsection{SMARTS: Interaction-Focused Multi-Agent Environment}

Whereas CARLA prioritizes visual fidelity, \citet{zhou2020smarts} designed
\emph{SMARTS} (Scalable Multi-Agent Reinforcement Learning Training School) with
the specific goal of \emph{diverse social vehicle behavior}. Its contributions
complement CARLA:
\begin{itemize}[noitemsep]
  \item A library of behavior models (IDM, MOBIL, aggressive, timid, distracted)
        enabling realistic, heterogeneous traffic.
  \item A gym-compatible interface with support for observations in multiple
        modalities (bird's-eye view, waypoints, agent observations).
  \item Scenario definition via SUMO road network files, enabling import of
        real-world road topologies.
\end{itemize}
SMARTS empirically demonstrates that training against diverse social vehicles
significantly improves the \emph{robustness} of learned policies, a result with
direct implications for sim-to-real transfer: real traffic is far more diverse
than static or homogeneous simulated opponents.

\subsection{NuPlan: Closed-Loop Planning Benchmark}

\citet{caesar2021nuplan} identify a systemic issue with the \emph{open-loop}
evaluation protocol dominant in prior benchmarks: when a planned trajectory
diverges from the ground-truth log trajectory, the system is penalized even if
the deviation is deliberate and safe (e.g., the ego-vehicle brakes to yield to an
aggressive merger). More critically, open-loop evaluation cannot assess whether
a planning algorithm handles \emph{interactive} scenarios correctly, since
surrounding agents do not react to the ego's plan.

NuPlan introduces the first large-scale \emph{closed-loop planning benchmark}
with reactive simulation:
\begin{itemize}[noitemsep]
  \item Over 1,200 hours of real-world driving log data across Las Vegas, Boston,
        Pittsburgh, and Singapore.
  \item Reactive agent models that respond to the ego-vehicle's maneuvers in
        simulation, providing a realistic assessment of interaction quality.
  \item A challenge suite with 1,080 unique scenarios spanning intersections,
        merge zones, and unprotected turns.
  \item Evaluation metrics including the \emph{Weighted Multi-metric Score} (WMS)
        combining progress, comfort, collision avoidance, and traffic rule
        adherence.
\end{itemize}
The NuPlan challenge has become the standard benchmark for comparing learned
planners (imitation-based, RL-based, and LLM-based) on an equal footing.

\subsection{Summary: Simulator--Benchmark Landscape}

Table~\ref{tab:simulators} summarizes the key properties of the surveyed platforms.
A clear trend is the shift from perception-focused datasets (nuScenes) to
interaction-and-planning-focused closed-loop benchmarks (NuPlan, SMARTS, MACAD),
reflecting the community's recognition that perception metrics do not predict
planning safety.

\begin{table}[htbp]
\centering
\caption{Comparison of simulation environments and benchmarks.}
\label{tab:simulators}
\begin{tabular}{lllll}
\toprule
\textbf{Platform} & \textbf{Loop} & \textbf{Multi-Agent} &
\textbf{Social Vehicles} & \textbf{Primary Use} \\
\midrule
CARLA \citep{dosovitskiy2017carlaopenurbandriving} & Both & Limited & Rule-based &
IL/RL training \& eval \\
SMARTS \citep{zhou2020smarts} & Closed & Full & Diverse models &
MARL training \\
MACAD-Gym \citep{palanisamy2020multi} & Closed & Full & CARLA agents &
Connected AD research \\
NuPlan \citep{caesar2021nuplan} & Closed & Full & Reactive &
Planning benchmark \\
\bottomrule
\end{tabular}
\end{table}

% ---------------------------------------------------------
\section{Open Challenges and Future Directions}
\label{sec:challenges}

This section synthesizes the residual limitations identified across all surveyed
paradigms into six structured open challenges.

\subsection{Formal Safety Verification for Learned Policies}
\label{sec:safe}

The RSS framework \citep{shalev2016safe} provides rule-based safety shields for
well-defined scenarios (following distance, right-of-way). However, LLM-based
cognitive agents generate decisions that are neither rule-based nor numerically
bounded. A GPT-4-generated decision can be linguistically coherent but physically
infeasible or safety-critical (``hallucinated'' decisions). Future work must
address \emph{formal verification of neuro-symbolic policies}: can RSS-like
invariants be proven to hold for the combined LLM--MPC architecture of
LanguageMPC \citep{sha2023languagempc}? Promising directions include conformal
prediction wrappers around LLM outputs and model-checking of LLM instruction
compliance.

\subsection{Scalable Cooperative Perception under Communication Constraints}
\label{sec:coop}

V2VNet \citep{wang2020v2vnet} demonstrates that intermediate feature sharing
outperforms both raw data and high-level track sharing. However, the work assumes
synchronized, lossless communication among vehicles within 70\,m. Real V2X
deployments face packet loss rates of 5--20\%, communication latency of 10--50\,ms,
and heterogeneous hardware clocks. Future work must develop \emph{asynchronous,
robust aggregation} methods that degrade gracefully under network impairments.
Furthermore, scaling cooperative perception to intersections with 20+ agents
requires hierarchical aggregation strategies beyond the pairwise GNN of V2VNet.

\subsection{Real-Time LLM Inference at the Edge}
\label{sec:latency}

LLM-based methods \citep{fu2023drive, wen2023dilu, mao2023gpt, sha2023languagempc}
rely on cloud-hosted or GPU-accelerated foundation models. At the time of those
publications (2023--2024), cloud API query latency for GPT-class models
typically ranged from hundreds of milliseconds to several seconds under standard
conditions---two orders of magnitude above the $\sim$10\,ms budget for real-time
control (while practical latency may improve in subsequent deployments). Even
LanguageMPC's 1--2\,Hz LLM cycle requires a stable low-latency network
connection. Edge deployment will
require: (a) model distillation \citep{kiran2021deep} to transfer reasoning
capabilities into compact on-board models, and (b) speculative decoding and
hardware-aware quantization to reduce latency while preserving reasoning quality.

\subsection{Generative World Models for Long-Tail Scenarios}
\label{sec:worldmodel}

A recurring theme across all paradigms is data scarcity in the tail of the driving
distribution. Perturbation-based data augmentation (ChauffeurNet \citep{bansal2018chauffeur})
and diverse social vehicle models (SMARTS \citep{zhou2020smarts}) partially address
this, but \emph{systematic} coverage of rare events (road debris, emergency vehicle
override, multi-party collision avoidance) requires generative approaches.
Diffusion-based world models that synthesize photo-realistic, dynamically
consistent driving scenes conditional on rare scenario descriptions are an
emerging frontier. These models would enable MAAD agents to be trained and evaluated
on effectively unlimited synthetic long-tail scenarios.

\subsection{Bidirectional Prediction--Planning at Scale}
\label{sec:gametheory}

PiP \citep{song2020pip} demonstrates that conditioning prediction on a small set
of $K$ ego-vehicle candidate trajectories improves joint forecasting accuracy. This
is a first-order approximation to the full \emph{extensive-form game} between the
ego-vehicle and $N$ surrounding agents, which has exponentially many joint
strategy profiles. Extending PiP to multi-agent scenes with many interacting
vehicles, and connecting it to the CTDE training paradigm of MADDPG
\citep{lowe2017multi}, remains an open research problem. Differentiable game
solvers and Monte Carlo tree search over the joint action space are plausible
directions.

\subsection{Sim-to-Real Transfer and Domain Adaptation}
\label{sec:simreal}

All surveyed training environments (CARLA, SMARTS, MACAD-Gym, NuPlan) exhibit a
visual and dynamic reality gap relative to real-world deployment. While perception
domain adaptation has been studied extensively, \emph{behavioral} sim-to-real
transfer---ensuring that policies learned against simulated social vehicles generalize
to real human drivers---is less understood. The diversity of social vehicle models in
SMARTS \citep{zhou2020smarts} is a positive step, but human driving behavior is
long-tailed in ways that parametric behavior models cannot fully capture.

% ---------------------------------------------------------
\section{Conclusion}
\label{sec:conclusion}

This survey has organized the multi-agent autonomous driving literature around a
three-layer taxonomy and traced a coherent technological progression. We summarize
the key insights at each layer.

\textbf{Single-Agent Foundations.}
End-to-end behavior cloning \citep{bojarski2016end} established that neural
networks can implicitly learn the visual features required for driving from raw
sensor data, but its vulnerability to covariate shift is fundamental and not
remediable by data scale alone. ChauffeurNet \citep{bansal2018chauffeur} showed
that \emph{synthesizing} adverse states during training is more sample-efficient
than collecting them. PPO \citep{schulman2017proximal} provides a stable policy
gradient algorithm that has become the default optimizer in MAAD training loops.

\textbf{Multi-Agent Interaction and Coordination.}
Representing agent interactions has evolved from local, grid-based social pooling
\citep{alahi2016social} to efficient, long-range vectorized GNNs
\citep{gao2020vectornet}. The decoupling between prediction and planning, long
taken for granted, has been shown by PiP \citep{song2020pip} to cause systematic
failures in interactive scenarios; bi-directional coupling is necessary. On the
policy learning side, CTDE frameworks \citep{lowe2017multi} resolve non-stationarity
by providing a globally consistent critic during training. Safety-constrained
hierarchical RL \citep{shalev2016safe} provides a complementary solution by
enforcing hard physical constraints that bound the space of learned policies.
Connected platforms \citep{palanisamy2020multi, zhou2020smarts} operationalize
these ideas at research scale. V2VNet \citep{wang2020v2vnet} extends cooperation
from policy to perception, demonstrating that intermediate feature sharing is the
optimal V2X communication strategy.

\textbf{LLM-Based Cognitive Agents.}
Large language models bring qualitatively new capabilities to autonomous driving:
chain-of-thought scene reasoning \citep{fu2023drive}, memory-augmented continual
adaptation \citep{wen2023dilu}, in-context motion planning \citep{mao2023gpt},
and real-time-safe neuro-symbolic control \citep{sha2023languagempc}. The common
thread is that LLMs provide a \emph{prior} over driving semantics that allows
generalization to long-tail scenarios without expensive retraining.

\textbf{The Convergence Ahead.}
The evidence across all three layers points toward a common architectural
direction: \emph{hybrid systems} that combine data-driven neural components for
pattern recognition and trajectory generation, structured symbolic components for
safety-critical constraint enforcement, and foundation model components for
semantic reasoning and continual adaptation. Realizing this vision requires
progress on all six open challenges identified in Section~\ref{sec:challenges}.
The field is well-positioned to make that progress: the mathematical foundations
(POSG, CTDE, RSS), the empirical platforms (CARLA, SMARTS, NuPlan), and the
cognitive tools (LLMs, chain-of-thought, retrieval-augmented generation) are all
in place.

% ---------------------------------------------------------
\bibliographystyle{plainnat}
\bibliography{references}

\end{document}
